\section*{الفصل \ref{ch:b2:climate-proxies} --- الأرشيفات والمؤشّرات المناخية البديلة}

\begin{solution}{exo:b2:climate-proxies:1}
(أ) $\delta = (27.5 \times (-30) + 1350 \times 0)/(1350 + 27.5) = -825/1377.5 = -0.60$ بالجزء في الألف.
(ب) يحمل المحيط عندئذ الماء كلّه: $-0.60$ بالجزء في الألف، أي أخفّ بمقدار \num{0.60} بالجزء في الألف ممّا هو
اليوم. (ج) $27.5/1377.5 = \qty{2.0}{\percent}$.
\end{solution}

\begin{solution}{exo:b2:climate-proxies:2}
(أ) $T = 16.1 - 4.76\,(2.50 - (0 - 0.27)) = 16.1 - 4.76 \times 2.77 = 16.1 - 13.19 =
\qty{2.9}{\celsius}$. (ب) $T = 16.1 - 4.76\,(2.50 - 0.73) = 16.1 - 8.43 = \qty{7.7}{\celsius}$: فالصدفة
نفسها تعطي درجة حرارة أعلى بنحو خمس درجات إذا كان الماء أثقل، ولهذا يجب
معرفة تركيب الماء. (ج) الصدفة أثقل بمقدار $2.71 - 2.50 = 0.21$ بالجزء في الألف،
فكان ماؤها أبرد بمقدار $4.76 \times 0.21 = \qty{1.0}{\celsius}$، أي عند \qty{1.9}{\celsius}.
\end{solution}

\begin{solution}{exo:b2:climate-proxies:3}
(أ) $0.4\,\mathrm{e}^{0.9} = \qty{0.98}{mmol/mol}$ و$0.4\,\mathrm{e}^{1.8} =
\qty{2.42}{mmol/mol}$. (ب) العامل $\mathrm{e}^{0.9} = 2.46$؛ ولكلّ درجة $\mathrm{e}^{0.09}
= 1.094$، أي نموّ قدره \qty{9.4}{\percent} لكلّ درجة. (ج) $T = \ln(\text{Mg/Ca}/0.4)/0.09$:
$\ln 5/0.09 = \qty{17.9}{\celsius}$ و$\ln 2.5/0.09 = \qty{10.2}{\celsius}$.
\end{solution}

\begin{solution}{exo:b2:climate-proxies:4}
(أ) $2668 - 2213 + 1 = 456$ سنة. (ب) قد تفقد شجرة واحدة حلقة محلّيًّا، أو تُظهر حلقة
لا يمكن قراءتها؛ ومطابقة أنماط الحلقات العريضة والضيّقة عبر أشجار كثيرة تكشف
الحلقات المفقودة وتثبّت سنة كلّ حلقة، والاتّفاق بين الأشجار يقيس
موثوقية التأريخ. (ج) عند خطّ الأشجار يحدّ عامل واحد، هو برد موسم النموّ،
الأشجار كلّها، فتشترك في النمط نفسه للسنوات الجيّدة والسيّئة؛ أمّا في الأراضي المنخفضة المعتدلة الرطبة
فلا يهيمن عامل واحد، وتختلف أنماط الأشجار، وقد يفشل \omterm{def:b2:climate-proxies:dendro}{التأريخ المتقاطع}.
\end{solution}

\begin{solution}{exo:b2:climate-proxies:5}
(أ) $125 \times 40/(3682.2 - 125) = 5000/3557.2 = 1.41$ بالجزء في الألف؛ و$134 \times 40/3548.2 =
5360/3548.2 = 1.51$ بالجزء في الألف. (ب) مع $H - h = 3552.2$: $3900/3552.2 = 1.10$، و$4550/3552.2 = 1.28$
و$6500/3552.2 = 1.83$ بالجزء في الألف. (ج) $\delta_{\mathrm{ice}} = -1.0 \times 3552.2/130 = -27.3$ بالجزء في
الألف، وهو أثقل من أيّ غطاء جليدي حالي (غرينلاند من $-30$ إلى $-35$، والغطاء الجليدي لأمريكا الشمالية
$-40$، وجليد القارة القطبية الجنوبية حتى $-50$). فموازنة الكتلة مع جليد واقعي وقياس ماء المسام
لا يتّفقان؛ والفجوة البالغة نحو نصف جزء في الألف شكّ حقيقي في الطريقة.
\end{solution}

\begin{solution}{exo:b2:climate-proxies:6}
(أ) $0.372\,\mathrm{e}^{0.0936T} = 0.4\,\mathrm{e}^{0.09T}$ يعطي $\mathrm{e}^{0.0036T} =
0.4/0.372 = 1.0753$، فيكون $T = \ln 1.0753/0.0036 = \qty{20.2}{\celsius}$. (ب) عند
\qty{2}{\celsius}: $0.4\,\mathrm{e}^{0.18} = 0.479$ و$0.372\,\mathrm{e}^{0.1872} = 0.449$، فتكون
النسبة الموفَّقة $0.449/0.479 = 0.937$، أي أدنى بنسبة \qty{6.3}{\percent}. (ج) $T = \ln(0.449/0.4)/0.09 =
\ln 1.1225/0.09 = \qty{1.3}{\celsius}$: فالقانون المتوقَّع يقرأ الصدفة أبرد بنحو
\qty{0.7}{\celsius} ممّا ينبغي. وفي الماء البارد يهمّ اختيار المعايرة بمقدار
درجة.
\end{solution}

\begin{solution}{exo:b2:climate-proxies:7}
(أ) من الهولوسين إلى GI-1، \qty{112.19}{m} من الجليد لما عدده 2989 سنة: $11\,219/2989 =
\qty{3.75}{cm}$ في السنة؛ ومن GI-16 إلى GI-17، \qty{17.89}{m} لما عدده 1160 سنة: $1789/1160 =
\qty{1.54}{cm}$ في السنة. (ب) الجليد
يجري: فكلّما دُفنت طبقة انتشرت جانبيًّا وترقّقت، والطبقات الأعمق والأقدم
ترقّقت مدّة أطول. (ج) حين تصبح الطبقات رقيقة جدًّا، سنتيمترًا أو أقلّ، لا تعود الأجهزة
قادرة على تمييز سنة من التي تليها، ويجب أن يستعمل \omterm{def:b2:climate-proxies:age-model}{نموذج العمر} نقاط ربط أو نماذج جريان
بدلًا من ذلك.
\end{solution}

\begin{solution}{exo:b2:climate-proxies:8}
(أ) كسر الفترة هو $(1550 - 1492.45)/(1604.64 - 1492.45) = 57.55/112.19 = 0.513$،
فيكون العمر $11\,703 + 0.513 \times 2989 = 11\,703 + 1533 = \num{13236}$ سنة b2k. (ب) عام 1950 يسبق عام 2000 بمقدار 50
سنة: \num{13186} سنة قبل 1950. (ج) $50 + 0.513 \times (93 - 50) = 72$ سنة.
\end{solution}

\begin{solution}{exo:b2:climate-proxies:9}
(أ) تسهم كلّ طبقة مشكوك فيها بنصف سنة: $2 \times 93 = 186$ و$2 \times 1287 = 2574$
طبقة مشكوكًا فيها. (ب) $\frac12\sqrt{186} = 6.8$ سنة و$\frac12\sqrt{2574} = 25.4$ سنة.
(ج) $93/6.8 = 14$ و$1287/25.4 = 51$. ونسبة $\frac12 n$ إلى $\frac12\sqrt{n}$ هي
$\sqrt{n}$: فهي تنمو بالجذر التربيعي لعدد الطبقات المشكوك فيها، فيزداد الاختيار بين
القاعدتين أهمّية مع العمق.
\end{solution}

\begin{solution}{exo:b2:climate-proxies:10}
(أ) على خطّ الماء الجوّي تتغيّر \mbox{\textdir TLT $\delta$D} ثماني مرّات بقدر تغيّر \mbox{\textdir TLT $\delta^{18}$O}، فيكون
الميل $8 \times 0.8 = 6.4$ بالجزء في الألف من \mbox{\textdir TLT $\delta$D} لكلّ درجة. والانخفاض $441.7 - 396.6 = 45.1$
بالجزء في الألف، ومن ثَمّ تبرّد قدره $45.1/6.4 = \qty{7.0}{\celsius}$. (ب) كان المحيط الجليدي أثقل بنحو
$8 \times 1.0 = 8$ بالجزء في الألف في \mbox{\textdir TLT $\delta$D}؛ فالثلج المتكوّن من محيط اليوم كان سيكون
أخفّ بالمقدار نفسه، $-441.7 - 8 = -449.7$ بالجزء في الألف. فيصبح الانخفاض $53.1$ بالجزء في الألف والتبرّد
$53.1/6.4 = \qty{8.3}{\celsius}$. (ج) $9.39 - 0.09 = \qty{9.30}{\celsius}$، فيكون $45.1/9.30 = 4.85$
بالجزء في الألف لكلّ درجة. والسلسلة المنشورة مصحَّحة أيضًا لتغيّر ارتفاع الغطاء
الجليدي، وعلاقتها بدرجة الحرارة مأخوذة من نماذج تحاكي نظائر الماء، لا من
الميل المكاني وحده.
\end{solution}

\begin{solution}{exo:b2:climate-proxies:11}
(أ) لوغاريتمات النسب $-0.528$، $-0.223$، $0.191$، $0.501$، $0.904$، $1.227$؛ ومتوسطها
$0.345$ ومتوسط درجة الحرارة 14. ومع $S_{TT} = \sum (T - 14)^2 = 280$ و$S_{T\ell} =
\sum (T - 14)(\ell - 0.345) = 24.93$، يكون $A = 24.93/280 = \qty{0.0890}{\per\celsius}$ و$\ln B =
0.345 - 0.0890 \times 14 = -0.901$، فيكون $B = \qty{0.406}{mmol/mol}$. (ب) البواقي هي
$+0.017$، $-0.034$، $+0.023$، $-0.023$، $+0.025$ و$-0.009$، أي بضعة في المئة من النسبة، بلا
اتّجاه؛ والتوفيق قريب من $0.4\,\mathrm{e}^{0.09T}$. (ج) $T = (\ln 1.30 + 0.901)/0.0890 = (0.262 +
0.901)/0.0890 = \qty{13.1}{\celsius}$. (د) $T = (\ln 0.45 + 0.901)/0.0890 = (-0.799 +
0.901)/0.0890 = \qty{1.2}{\celsius}$. وتقع هذه القيمة دون أبرد عيّنة في المعايرة،
\qty{4}{\celsius}: فالقانون مُستقرَأ، ومعروف أن نسبة Mg/Ca في الأصداف تحيد عند درجات الحرارة المنخفضة
عن القانون الأسّي بتأثير أيون الكربونات.
\end{solution}

\begin{solution}{exo:b2:climate-proxies:12}
(أ) لكلّ وحدة من مساحة المحيط، يتوزّع ماء المحيط الحالي، $H$ عند 0، بين المخازن
$h_i$ عند $\delta_i$ والمحيط الباقي $H - h$، مع $h = \sum h_i$: $0 = \sum h_i\delta_i +
(H - h)\,\Delta\delta$. وبعد المرحلة الأولى $\Delta\delta = 60 \times 30/(3682.2 - 60) =
1800/3622.2 = 0.50$ بالجزء في الألف؛ وبعد الثانية $(1800 + 70 \times 45)/(3682.2 - 130) =
4950/3552.2 = 1.39$ بالجزء في الألف. (ب) $-4950/130 = -38.1$ بالجزء في الألف. (ج) القاعدة $1.39/130 = 0.0107$
بالجزء في الألف لكلّ متر، وتعطي $0.50/0.0107 = \qty{47}{m}$ للمرحلة الأولى بدل
\qty{60}{m}: فالجليد الأوّل كان أثقل من المتوسط، فكان كلّ متر منه يزيح المحيط
أقلّ. والقاعدة الخطّية المعايَرة على حالة واحدة تخطئ قراءة الحالات الأخرى حين يتغيّر متوسط تركيب
الجليد.
\end{solution}

\begin{solution}{pb:b2:climate-proxies:1}
\textbf{1.} $361.84 \times 10^6\ \unit{km^2} \times \qty{0.130}{km} = \qty{4.70e7}{km^3}$، أي نحو 47
مليون كيلومتر مكعّب.
\textbf{2.} $130 \times 40/(3682.2 - 130) = 5200/3552.2 = 1.46$ بالجزء في الألف.
\textbf{3.} $5000/3557.2 = 1.41$ و$5360/3548.2 = 1.51$ بالجزء في الألف.
\textbf{4.} $4550/3552.2 = 1.28$ و$6500/3552.2 = 1.83$ بالجزء في الألف. فتركيب الجليد
يحرّك الجواب بمقدار \num{0.55} بالجزء في الألف، ومستوى البحر بمقدار \num{0.10}: فالجليد أهمّ.
\textbf{5.} $\delta_{\mathrm{ice}} = -1.0 \times 3552.2/130 = -27.3$ بالجزء في الألف.
\textbf{6.} $40/3682.2 = 0.0109$ بالجزء في الألف لكلّ متر، وتقريبه هو القيمة \num{0.011} المستمدّة من
مستويات البحر المؤرَّخة.
\textbf{7.} مع ثبات $\delta^{18}\mathrm{O}_{\mathrm{sw}}$، يكون $\mathrm{d}T = -4.76\,\mathrm{d}\delta^{18}
\mathrm{O}_{\mathrm{b}}$، فإن $\mathrm{d}T = -1$ يعطي $\mathrm{d}\delta^{18}\mathrm{O}_{\mathrm{b}} =
1/4.76 = 0.21$ بالجزء في الألف.
\textbf{8.} $\delta^{18}\mathrm{O}_{\mathrm{b}} = (16.1 - T)/4.76 + \delta^{18}\mathrm{O}_{\mathrm{sw}}
- 0.27 = 14.1/4.76 - 0.27 = 2.96 - 0.27 = 2.69$ بالجزء في الألف.
\textbf{9.} $17.1/4.76 + 1.0 - 0.27 = 3.59 + 0.73 = 4.32$ بالجزء في الألف. والتغيّر $4.32 - 2.69 =
1.63$ بالجزء في الألف: منها \num{1.00} من الماء و$3.0 \times 0.21 = 0.63$ من التبرّد بمقدار
\qty{3}{\celsius}.
\textbf{10.} $T = \ln(2.0/0.4)/0.09 = \ln 5/0.09 = \qty{17.9}{\celsius}$.
\textbf{11.} $0.18/0.1 = \qty{1.8}{\celsius}$ من التبرّد.
\textbf{12.} نسبة Mg/Ca نسبة بين عنصرين في الكالسيت، تحدّدها درجة حرارة
التكلّس؛ ولا علاقة لها بنظائر الأكسجين، فهي لا تتأثّر بتركيب
الماء الذي يغيّره الجليد.
\textbf{13.} $4.92 - 3.23 = 1.69$ بالجزء في الألف.
\textbf{14.} بحلّ $\Delta = h \times 40/(H - h)$ بالنسبة إلى $h$: $h = H\Delta/(\Delta + 40) = 3682.2
\times 1.69/41.69 = \qty{149}{m}$، أي أكثر من 125 إلى \qty{134}{m} التي تشهد بها شواهد مستوى البحر:
فجزء من التغيّر لا بدّ أن يكون تبرّدًا.
\textbf{15.} $1.69 - 1.46 = 0.23$ بالجزء في الألف، و$4.76 \times 0.23 = \qty{1.1}{\celsius}$.
\textbf{16.} $1.69 - 1.00 = 0.69$ بالجزء في الألف، و$4.76 \times 0.69 = \qty{3.3}{\celsius}$ (3.28).
\textbf{17.} $0.23/0.25 = \qty{0.9}{\celsius}$ و$0.69/0.25 = \qty{2.8}{\celsius}$.
\textbf{18.} تبرّد ماء شمال الأطلسي العميق من 2--\qty{4}{\celsius} إلى \qty{-1.1}{\celsius}،
أي بمقدار 3.1 إلى \qty{5.1}{\celsius}؛ وماء القاع القطبي الجنوبي من 0--\qty{1}{\celsius} إلى
\qty{-1.9}{\celsius}، أي بمقدار 1.9 إلى \qty{2.9}{\celsius}؛ وتعطي نسبة Mg/Ca في شمال غرب الأطلسي
\qty{1.8}{\celsius}. وكلّها تقع بين نحو 2 و\qty{5}{\celsius}، حول 2.8 إلى
\qty{3.3}{\celsius} في تصحيح ماء المسام، وفوق \qty{1}{\celsius} في موازنة
الكتلة مع جليد عند $-40$ بالجزء في الألف.
\textbf{19.} $-396.6 - (-441.7) = 45.1$ بالجزء في الألف.
\textbf{20.} $8 \times 0.8 = 6.4$ بالجزء في الألف لكلّ درجة، فيكون $45.1/6.4 = \qty{7.0}{\celsius}$.
\textbf{21.} $9.39 - 0.09 = \qty{9.30}{\celsius}$، و$45.1/9.30 = 4.85$ بالجزء في الألف لكلّ درجة.
\textbf{22.} القبّة C، \qty{9.3}{\celsius}، أكثر من السطح العالمي، 5 إلى \qty{7}{\celsius}، وهذا أكثر
من المحيط العميق، نحو \qty{3}{\celsius}.
\textbf{23.} الماء العميق بارد أصلًا اليوم، من 0 إلى \qty{4}{\celsius}، ولا يمكن لماء البحر أن
يبرد كثيرًا دون نقطة تجمّده، التي كاد ماء القاع الجليدي يبلغها عند
\qty{-1.9}{\celsius}: فلم يكن للمحيط العميق إلا بضع درجات يفقدها، بينما ليس للهواء فوق القارة القطبية الجنوبية
مثل هذا الحدّ الأدنى.
\textbf{24.} مع تصحيح ماء المسام، كان المحيط العميق \textbf{أبرد بنحو \qty{3.3}{\celsius}}
في الذروة الجليدية الأخيرة منه اليوم، قريبًا من 2 إلى \qty{5}{\celsius} في
تقديرات الماء العميق المستقلّة.
\end{solution}
